\documentclass[11pt, a4paper]{article}

% --- CODIFICA, LINGUA E TIPOGRAFIA ---
\usepackage[utf8]{inputenc}
\usepackage[T1]{fontenc}
\usepackage[italian]{babel}
\usepackage{lmodern}
\usepackage{microtype} 

% --- GEOMETRIA E LAYOUT ---
\usepackage[top=2.5cm, bottom=2.5cm, left=2.5cm, right=2.5cm, headheight=15pt]{geometry}
\usepackage{fancyhdr}
\usepackage{tabularx}
\usepackage{booktabs}
\usepackage{graphicx}
\usepackage{caption}
\usepackage{xcolor}

% --- COLLEGAMENTI E METADATI PDF ---
\usepackage[
    colorlinks=true,
    linkcolor=black,
    urlcolor=blue!70!black,
    citecolor=black,
    pdftitle={Verbale Interno N. 5 - 13oolean},
    pdfauthor={13oolean}
]{hyperref}

% --- DEFINIZIONE METADATI E COMANDI RIUTILIZZABILI ---
\newcommand{\GroupName}{13oolean}
\newcommand{\GroupEmail}{\href{mailto:gruppo13oolean@gmail.com}{gruppo13oolean@gmail.com}}
\newcommand{\DocTitle}{Verbale Riunione Interna \#5}
\newcommand{\DocDate}{2026-11-04}

% --- CONFIGURAZIONE INTESTAZIONI E PIÈ DI PAGINA ---
\pagestyle{fancy}
\fancyhf{}
\lhead{\textbf{\GroupName}}
\chead{\DocTitle}
\rhead{\DocDate}
\cfoot{\thepage}
\renewcommand{\headrulewidth}{0.4pt}
\renewcommand{\footrulewidth}{0pt}

% ==============================================================================
\begin{document}

% ------------------------------------------------------------------------------
% FRONTESPIZIO
% ------------------------------------------------------------------------------
\begin{titlepage}
    \centering
    
    \includegraphics[width=0.45\textwidth]{../../../shared/logo_unipd.png}\par
    \vspace{0.8cm}
    
    {\scshape\large Università degli Studi di Padova\par}
    {\large Corso di Laurea Triennale in Informatica\par}
    {\large Corso di Ingegneria del Software -- A.A. 2026/2027\par}
    
    \vspace{1cm}
    \includegraphics[width=0.35\textwidth]{../../../shared/logo_gruppo.jpeg}\par
    \vspace{0.8cm}

    {\large Gruppo: \textbf{\GroupName}\par}
    {\large Email: \GroupEmail\par}

    \vfill

    {\Huge\bfseries \DocTitle\par}
    \vspace{0.3cm}
    {\large Data: \DocDate\par}
    
    \vspace{1.5cm}
\end{titlepage}

% ------------------------------------------------------------------------------
% INDICE E INIZIO CONTENUTO
% ------------------------------------------------------------------------------
\pagenumbering{Roman} % Numerazione romana per le sezioni iniziali
\tableofcontents
\clearpage

\pagenumbering{arabic} % Numerazione araba per il contenuto effettivo

\section{Informazioni Generali}
\begin{itemize}
    \item \textbf{Luogo / Piattaforma:} Server Discord
    \item \textbf{Data:} \DocDate
    \item \textbf{Orario:} 15:30 -- 17:50
    \item \textbf{Responsabile:} Michele Tosato
    \item \textbf{Verificatore:} Andrea Guerra
    \item \textbf{Scriba:} Davide Menegaldo
    \item \textbf{Partecipanti:}
        \begin{itemize}
            \item Andrea Guerra
            \item Francesco Cirigliano
            \item Gabriele Maria D'Amico
            \item Giacomo Lazzarotto
            \item Davide Menegaldo
            \item Michele Tosato
        \end{itemize}
\end{itemize}
\clearpage


\section{Ordine del Giorno}
\begin{enumerate}
    \item Analisi dell'esito della candidatura e attuazione delle osservazioni del docente.
    \item Definizione dei criteri di accettazione per le Pull Request e standardizzazione dei commit.
    \item Configurazione delle label per le issue su GitHub.
    \item Organizzazione degli strumenti di gestione per la rotazione ruoli (Google Sheets) e lo scadenziario (Google Calendar).
    \item Pianificazione delle attività verso la baseline RTB: Glossario, Analisi dei Requisiti e Norme di Progetto.
\end{enumerate}


\section{Resoconto della Riunione}

\subsection{Attuazioni a seguito dell'esito della candidatura}
A seguito dell'assegnazione ufficiale del capitolato C0, il gruppo ha preso in esame i riscontri forniti dal docente e ha definito i seguenti interventi correttivi:
\begin{itemize}
    \item Aggiornamento del logo dell'Università di Padova con il nuovo sigillo istituzionale su tutti i documenti e aggiornamento del template documentale;
    \item Revisione dei verbali pregressi, inserimento di un solo verificatore, correzione del versionamento e aggiornamento dei nomi dei file;
    \item Impostazione del formato data \texttt{AAAA-MM-DD} su tutti i documenti di progetto;
    \item Riorganizzazione della home page del sito web con apertura dei documenti in una nuova scheda e introduzione della sezione per i documenti in evidenza.
\end{itemize}

\subsection{Definizione criteri di accettazione delle PR e standardizzazione dei commit}
Al fine di consolidare il flusso di lavoro sul repository, sono stati formalizzati i criteri di accettazione per la chiusura delle Pull Request come segue:
\begin{itemize}
    \item Il documento è stato redatto su un branch dedicato, separato da \texttt{main};
    \item Il testo è stato riletto ed è privo di refusi;
    \item Struttura, titoli e formattazione seguono il template e le convenzioni di stile del progetto;
    \item Link, riferimenti, immagini e tabelle sono corretti e funzionanti;
    \item Il changelog (o registro delle modifiche) riporta le variazioni apportate;
    \item Il documento viene generato correttamente (build PDF o Markdown) senza errori o warning;
    \item È stata aperta una Pull Request con una breve descrizione dei contenuti modificati;
    \item Almeno un revisore ha verificato il documento e approvato la Pull Request.
\end{itemize}
È stata inoltre concordata la struttura standardizzata dei commit e stabilita la cancellazione dei branch di lavoro a seguito del merge.

\subsection{Configurazione delle label per le issue}
Per categorizzare le issue nel repository GitHub, il gruppo ha deciso di introdurre le label suddivise nella seguente maniera: \texttt{Priority:\ High}, \texttt{Medium}, \texttt{Low}; \texttt{ Type:\ Code}, \texttt{Documentation}, \texttt{Web}; \texttt{ Fix}.

\subsection{Strumenti di gestione: rotazione ruoli e scadenziario}
Per pianificare e monitorare la rotazione dei ruoli e la rendicontazione oraria individuale, il gruppo ha approvato l'adozione di un foglio di calcolo condiviso su Google Sheets. È stato inoltre integrato Google Calendar per la gestione delle scadenze e delle pianificazioni. Il gruppo si riserva di valutare anche altri strumenti e un'eventuale futura integrazione.

\subsection{Pianificazione delle attività verso la baseline RTB}
In vista della preparazione della milestone RTB, il gruppo ha definito i compiti propedeutici per i documenti cardine:
\begin{itemize}
    \item Redazione del Glossario;
    \item Stesura del documento e formalizzazione delle metodologie di lavoro per le Norme di Progetto;
    \item Studio individuale approfondito del capitolato C0 per redigere una bozza iniziale dell'Analisi dei Requisiti, compito affidato a tutti i componenti in vista della riunione del 2026-11-09.
\end{itemize}


\clearpage
\section{Decisioni Prese}

\begin{table}[htbp]
\centering
\caption{Riepilogo delle decisioni assunte durante la riunione}
\label{tab:decisioni}
\begin{tabularx}{\textwidth}{l X}
\toprule
\textbf{Codice} & \textbf{Decisione} \\
\midrule
DEC-5.1 & Stesura del verbale della riunione interna del 2026-11-04. \\
DEC-5.2 & Redazione e aggiornamento del Diario di Bordo. \\
DEC-5.3 & Introduzione delle label per il tracciamento e la priorità delle issue su GitHub. \\
DEC-5.4 & Revisione e correzione dei verbali pregressi. \\
DEC-5.5 & Riorganizzazione della home page del sito web e adozione del formato data \texttt{AAAA-MM-DD}. \\
DEC-5.6 & Stesura del verbale della riunione interna del 2026-10-28. \\
DEC-5.7 & Redazione del Glossario. \\
DEC-5.8 & Stesura del template e delle metodologie di lavoro per le Norme di Progetto. \\
DEC-5.9 & Stesura di una prima bozza dell'Analisi dei Requisiti. \\ 
DEC-5.10 & Approvazione dei criteri di accettazione per le Pull Request. \\
DEC-5.11 & Adozione di una struttura standardizzata per i commit e chiusura dei branch post-merge. \\
DEC-5.12 & Adozione di Google Sheets per la rotazione dei ruoli e di Google Calendar per lo scadenziario. \\
\bottomrule
\end{tabularx}
\end{table}

\section{Azioni da Svolgere}

\begin{table}[htbp]
\centering
\caption{Elenco delle azioni e task assegnati}
\label{tab:azioni}
\begin{tabularx}{\textwidth}{l l X l c}
\toprule
\textbf{Codice} & \textbf{Rif.} & \textbf{Descrizione Task} & \textbf{Assegnatario} & \textbf{Scadenza} \\
\midrule
T-5.1 & DEC-5.1 & Stesura verbale riunione 2026-11-04 & Davide & 2026-11-06 \\
T-5.2 & DEC-5.2 & Redazione del Diario di Bordo & Michele & 2026-11-06 \\
T-5.3 & DEC-5.3 & Aggiunta delle label alle issue su GitHub & Andrea & 2026-11-06 \\
T-5.4 & DEC-5.4 & Correzione dei verbali pregressi & Francesco & 2026-11-06 \\
T-5.5 & DEC-5.5 & Riorganizzazione home page del sito e formato date & Gabriele & 2026-11-06 \\
T-5.6 & DEC-5.6 & Stesura verbale riunione 2026-10-28 & Giacomo & 2026-11-06 \\
T-5.7 & DEC-5.7 & Redazione del Glossario & Giacomo & 2026-11-09 \\
T-5.8 & DEC-5.8 & Stesura template metodologie di lavoro & Francesco & 2026-11-09 \\
T-5.9 & DEC-5.9 & Redazione di una bozza dell'Analisi dei Requisiti & Tutti & 2026-11-09 \\
\bottomrule
\end{tabularx}
\end{table}



\end{document}
